% ============================================================================
%  ECE 431 -- Sheet 4 Solutions
%  Style: sheet-answering-style.md
% ============================================================================

\begin{question}[Question 1]The transmittance of a symmetrical Fabry-Perot resonator was measured by
using light from a tunable monochromatic light source. The transmittance versus frequency
exhibits periodic pulses of period 150\,MHz, each of width (FWHM) 5\,MHz. Assuming that the
medium within the resonator mirrors is a gas with $n=1$, determine the length and finesse of
the resonator. Assuming that the only source of loss is associated with the mirrors, find their
reflectances.
\end{question}

\bigskip
\textbf{\large Solution --- Question 1}

\medskip
\textbf{Resonator length and finesse.}

The free spectral range (FSR) is the period of the periodic pulses, so $\nu_F = 150$\,MHz. The FWHM linewidth is $\delta\nu = 5$\,MHz.
The resonator length $L$ is derived directly from the FSR for an empty cavity ($n=1$):

\begin{equation*}
  \nu_F = \frac{c}{2L} \quad\Longrightarrow\quad L = \frac{c}{2\nu_F} = \frac{3\times10^{8}}{2(150\times10^{6})}
\end{equation*}

\begin{equation*}
  \boxed{L = 1\;\text{m.}}
\end{equation*}

The finesse $F$ characterizes the sharpness of the resonances, given by the ratio of the FSR to the linewidth:

\begin{equation*}
  F = \frac{\nu_F}{\delta\nu} = \frac{150\,\text{MHz}}{5\,\text{MHz}}
\end{equation*}

\begin{equation*}
  \boxed{F = 30.}
\end{equation*}

\textbf{Mirror reflectances.}

Assuming identical mirrors ($R_1 = R_2 = R$) and no internal loss ($\alpha_s=0$), the finesse is related to the reflectance by:

\begin{equation*}
  F \approx \frac{\pi\sqrt{R}}{1-R} \quad\Longrightarrow\quad 30 = \frac{\pi\sqrt{R}}{1-R}
\end{equation*}

This leads to a quadratic equation in $\sqrt{R}$:
$30 - 30R - \pi\sqrt{R} = 0$.
Solving for $\sqrt{R}$ via the quadratic formula:

\begin{equation*}
  \sqrt{R} = \frac{-\pi + \sqrt{\pi^2 - 4(-30)(30)}}{2(-30)} = \frac{\pi - \sqrt{\pi^2 + 3600}}{-60}
\end{equation*}

Choosing the physically meaningful positive root yields $\sqrt{R} \approx 0.949$.

\begin{equation*}
  \boxed{R \approx 0.90.}
\end{equation*}

Thus, each mirror has a reflectance of approximately 90\%.

% ============================================================================

\begin{question}[Question 2]What time does it take for the optical energy stored in a resonator of
finesse $F=100$, length $L=50$\,cm, and refractive index $n=1$, to decay to one-half of its
initial value?
\end{question}

\bigskip
\textbf{\large Solution --- Question 2}

\medskip
\textbf{Photon lifetime.}

The photon lifetime $\tau_p$ is the $e$-folding decay time of the optical energy inside the cavity. It is inversely proportional to the resonator linewidth $\delta\nu$:

\begin{equation*}
  \tau_p = \frac{1}{2\pi\delta\nu}
\end{equation*}

The linewidth can be found from the given finesse and free spectral range ($\nu_F = c/2L$):

\begin{equation*}
  \delta\nu = \frac{\nu_F}{F} = \frac{c}{2LF} = \frac{3\times10^{8}}{2(0.5)(100)} = 3\times10^{6}\;\text{Hz} = 3\,\text{MHz}.
\end{equation*}

Therefore, the photon lifetime is:

\begin{equation*}
  \tau_p = \frac{1}{2\pi(3\times10^{6})} \approx 53.05\,\text{ns}.
\end{equation*}

\textbf{Half-life of stored energy.}

The optical energy $W(t)$ decays exponentially according to $W(t) = W(0) e^{-t/\tau_p}$. We need the time $t_{1/2}$ when $W(t_{1/2}) = \frac{1}{2}W(0)$:

\begin{equation*}
  \frac{1}{2} = e^{-t_{1/2}/\tau_p} \quad\Longrightarrow\quad t_{1/2} = \tau_p \ln(2)
\end{equation*}

\begin{equation*}
  \boxed{t_{1/2} = (53.05\,\text{ns})\ln(2) \approx 36.77\,\text{ns}.}
\end{equation*}

% ============================================================================

\begin{question}[Question 3]\begin{enumerate}[label=(\alph*)]
  \item Determine the spacing between adjacent resonance frequencies in a resonator constructed
        of two parallel planar mirrors separated by a distance $L=15$\,cm in air ($n=1$).
  \item A transparent plate of thickness $L=2.5$\,cm and refractive index $n=1.5$ is placed
        inside the resonator and is tilted slightly to prevent light reflected from the plate
        from reaching the mirrors. Determine the spacing between the resonance frequencies of
        the resonator.
\end{enumerate}
\end{question}

\bigskip
\textbf{\large Solution --- Question 3}

\medskip
\textbf{Part (a): Empty resonator spacing.}

The spacing between adjacent resonance frequencies is the free spectral range (FSR). For a planar mirror cavity of length $L=15$\,cm in air ($n=1$):

\begin{equation*}
  \nu_F = \frac{c}{2L} = \frac{3\times10^{8}}{2(0.15)}
\end{equation*}

\begin{equation*}
  \boxed{\nu_F = 10^{9}\;\text{Hz} = 1\;\text{GHz}.}
\end{equation*}

\textbf{Part (b): Plate-loaded resonator spacing.}

When a transparent plate of thickness $L_p = 2.5$\,cm and index $n_p = 1.5$ is inserted, the physical length remains $L=15$\,cm, but the optical path length changes. The plate is tilted only slightly, so the path length through the plate is approximately $L_p$. The effective optical length $L_{\text{opt}}$ is the sum of the air path and the plate path:

\begin{equation*}
  L_{\text{opt}} = (L - L_p)n_{\text{air}} + L_p n_p = (15 - 2.5)(1) + (2.5)(1.5) = 12.5 + 3.75 = 16.25\,\text{cm}.
\end{equation*}

The new resonance spacing is determined by this extended optical length:

\begin{equation*}
  \nu_F' = \frac{c}{2L_{\text{opt}}} = \frac{3\times10^{8}}{2(0.1625)}
\end{equation*}

\begin{equation*}
  \boxed{\nu_F' \approx 923.08\;\text{MHz}.}
\end{equation*}

The insertion of the higher-index plate effectively "lengthens" the cavity, slightly decreasing the free spectral range.

% ============================================================================

\begin{question}[Question 4]Derive an expression for the resonance frequencies of the three-mirror
cavity ring resonator shown in the figure. Assume that each mirror reflection introduces a phase
shift of $\pi$.
\end{question}

\bigskip
\textbf{\large Solution --- Question 4}

\medskip
\textbf{Ring resonator phase condition.}

For a wave circulating in a three-mirror ring resonator with total perimeter $L$, resonance occurs when the accumulated phase after one complete round trip is an integer multiple of $2\pi$. The total phase includes the propagation phase $kL$ and the discrete phase shifts from the reflections.
With three mirrors, each contributing a $\pi$ phase shift, the total reflection phase shift is $3\pi$. The resonance condition is:

\begin{equation*}
  kL - 3\pi = 2\pi q
\end{equation*}

where $q$ is an integer and $k = 2\pi\nu/c$. Substituting $k$:

\begin{equation*}
  2\pi\frac{\nu}{c}L = 2\pi q + 3\pi \quad\Longrightarrow\quad \frac{2\nu L}{c} = 2q + 3.
\end{equation*}

Solving for the resonance frequencies $\nu_q$:

\begin{equation*}
  \boxed{\nu_q = \left(q + \frac{3}{2}\right)\frac{c}{L}.}
\end{equation*}

Unlike a linear Fabry-Perot cavity where the FSR is $c/2L$, the ring resonator's FSR is $c/L$ because the wave travels the entire perimeter in exactly one round trip. Furthermore, the odd number of $\pi$ phase shifts offsets the resonance comb by exactly half an FSR from the zero-frequency origin.

% ============================================================================

\begin{question}[Question 5]A Doppler-broadened gas laser operates at 515\,nm in a symmetric resonator
with two mirrors each of power reflectivity 0.9 separated by a distance of 50\,cm. The spectral
window within which lasing occurs is of bandwidth $B=1.5$\,GHz. The refractive index $n=1$.
Calculate
\begin{enumerate}[label=(\arabic*)]
  \item the resonator loss coefficient,
  \item the photon lifetime,
  \item the number of longitudinal modes inside the laser cavity,
  \item the finesse of the laser cavity,
  \item the spectral width of each longitudinal mode.
\end{enumerate}
\end{question}

\bigskip
\textbf{\large Solution --- Question 5}

\medskip
\textbf{Part (1): Resonator loss coefficient.}

Assuming the only source of loss is the transmission through the two partially reflecting mirrors, the distributed loss coefficient is:

\begin{equation*}
  \alpha_r = \frac{1}{2L}\ln\left(\frac{1}{R_1 R_2}\right) = \frac{1}{L}\ln\left(\frac{1}{R}\right) = \frac{1}{0.5}\ln\left(\frac{1}{0.9}\right)
\end{equation*}

\begin{equation*}
  \boxed{\alpha_r \approx 0.211\;\text{m}^{-1} = 2.11\times10^{-3}\;\text{cm}^{-1}.}
\end{equation*}

\textbf{Part (2): Photon lifetime.}

The photon lifetime dictates how long energy remains trapped within the cavity before being lost. Using the loss coefficient:

\begin{equation*}
  \tau_p = \frac{1}{c\alpha_r} = \frac{1}{(3\times10^{8})(0.2107)}
\end{equation*}

\begin{equation*}
  \boxed{\tau_p \approx 1.58\times10^{-8}\;\text{s} = 15.8\;\text{ns}.}
\end{equation*}

\textbf{Part (3): Number of longitudinal modes.}

The free spectral range (FSR) is the spacing between adjacent modes:

\begin{equation*}
  \nu_F = \frac{c}{2L} = \frac{3\times10^{8}}{2(0.5)} = 300\;\text{MHz}.
\end{equation*}

The number of oscillating modes $M$ is the spectral window $B$ divided by the modal spacing:

\begin{equation*}
  M \approx \frac{B}{\nu_F} = \frac{1.5\times10^{9}}{300\times10^{6}}
\end{equation*}

\begin{equation*}
  \boxed{M = 5\;\text{modes}.}
\end{equation*}

\textbf{Part (4): Finesse of the cavity.}

For symmetric mirrors with identical reflectivity $R=0.9$:

\begin{equation*}
  F \approx \frac{\pi\sqrt{R}}{1 - R} = \frac{\pi\sqrt{0.9}}{1 - 0.9}
\end{equation*}

\begin{equation*}
  \boxed{F \approx 29.8.}
\end{equation*}

\textbf{Part (5): Spectral width of each mode.}

The passive FWHM linewidth of a single longitudinal mode is the FSR divided by the finesse:

\begin{equation*}
  \delta\nu = \frac{\nu_F}{F} = \frac{300\,\text{MHz}}{29.8}
\end{equation*}

\begin{equation*}
  \boxed{\delta\nu \approx 10.1\;\text{MHz}.}
\end{equation*}

% ============================================================================

\begin{question}[Question 6]A 632.8\,nm HeNe laser has the following characteristics: (1) a resonator
with 97\% and 100\% mirror power reflectances and negligible internal loss; (2) a
Doppler-broadened atomic transition with Doppler linewidth $\Delta\nu_{D}=1.5$\,GHz; and (3) a
small-signal peak gain coefficient $\gamma_{o}(\nu_{o})=2.5\times10^{-3}$\,cm$^{-1}$. The
refractive index $n=1$.
\begin{enumerate}[label=(\alph*)]
  \item Find the maximum cavity length which ensures single longitudinal mode operation,
        assuming that the longitudinal mode frequency coincides with the 632.8\,nm laser line.
        Then, use this cavity length to find
  \item the photon flux density $\Phi$ inside the cavity and the output laser intensity when
        the saturation photon-flux density at the centre wavelength is
        $1.4\times10^{19}$\,photons/cm$^{2}$-s, and
  \item the FWHM linewidth $\delta\nu$ of the single longitudinal laser mode.
\end{enumerate}
\end{question}

\bigskip
\textbf{\large Solution --- Question 6}

\medskip
\textbf{Part (a): Maximum single-mode cavity length.}

To enforce single longitudinal mode operation when the center mode is perfectly aligned with the gain peak ($\nu_o$), we must guarantee that the adjacent modes located at $\nu_o \pm \nu_F$ have a small-signal gain strictly less than the resonator loss ($\gamma_o(\nu_o \pm \nu_F) < \alpha_r$).

First, we determine the cavity loss $\alpha_r$ as a function of the cavity length $L$ (in cm):

\begin{equation*}
  \alpha_r = \frac{1}{2L}\ln\left(\frac{1}{R_1 R_2}\right) = \frac{1}{2L}\ln\left(\frac{1}{0.97}\right) \approx \frac{0.01523}{L}.
\end{equation*}

The inhomogeneously broadened Gaussian gain profile yields the gain at the adjacent modes:

\begin{equation*}
  \gamma_o(\nu_o \pm \nu_F) = \gamma_o(\nu_o) \exp\left[ -4\ln 2 \left( \frac{\nu_F}{\Delta\nu_D} \right)^2 \right].
\end{equation*}

Using $\nu_F = c/2L = 15\times10^{9}/L$ Hz (with $L$ in cm) and $\Delta\nu_D = 1.5\times10^{9}$ Hz, the ratio becomes $\nu_F/\Delta\nu_D = 10/L$. The critical length occurs when gain equals loss:

\begin{equation*}
  2.5\times10^{-3} \exp\left[ -4\ln 2 \left(\frac{10}{L}\right)^2 \right] = \frac{0.01523}{L}
  \quad\Longrightarrow\quad 2.5 \exp\left( -\frac{277.26}{L^2} \right) = \frac{15.23}{L}.
\end{equation*}

Solving this transcendental equation numerically yields the critical length $L_{\max}$:

\begin{equation*}
  \boxed{L_{\max} \approx 16.6\;\text{cm}.}
\end{equation*}
Any length shorter than this will push the adjacent modes further out, guaranteeing their gain is below the loss threshold.

\textbf{Part (b): Cavity flux and output intensity.}

For an inhomogeneously broadened medium, the steady-state saturated gain satisfies:

\begin{equation*}
  \gamma(\nu_o) = \frac{\gamma_o(\nu_o)}{\sqrt{1 + \Phi/\Phi_s}} = \alpha_r.
\end{equation*}

Using $L = 16.63$\,cm, the threshold loss is $\alpha_r = 0.01523 / 16.63 \approx 0.916\times10^{-3}$\,cm$^{-1}$. Solving for the internal flux $\Phi$:

\begin{equation*}
  \sqrt{1 + \frac{\Phi}{\Phi_s}} = \frac{2.5\times10^{-3}}{0.916\times10^{-3}} \approx 2.73
  \quad\Longrightarrow\quad \frac{\Phi}{\Phi_s} = 2.73^2 - 1 \approx 6.45.
\end{equation*}

\begin{equation*}
  \boxed{\Phi = 6.45(1.4\times10^{19}) \approx 9.03\times10^{19}\;\text{photons/cm}^{2}\text{-s}.}
\end{equation*}

The output flux emerges from the partially transmitting mirror ($T_1 = 1 - 0.97 = 0.03$). Given $\Phi$ represents the total two-way circulating flux, the forward-traveling component is $\Phi/2$:

\begin{equation*}
  \Phi_{\text{out}} = \frac{T_1}{2}\Phi = \frac{0.03}{2}(9.03\times10^{19}) \approx 1.35\times10^{18}\;\text{photons/cm}^{2}\text{-s}.
\end{equation*}

The output intensity is $I_{\text{out}} = h\nu\Phi_{\text{out}}$, with photon energy $h\nu = hc/\lambda \approx 3.14\times10^{-19}$\,J:

\begin{equation*}
  \boxed{I_{\text{out}} \approx (3.14\times10^{-19})(1.35\times10^{18}) \approx 0.424\;\text{W/cm}^{2} = 424\;\text{mW/cm}^{2}.}
\end{equation*}

\textbf{Part (c): FWHM linewidth.}

The passive cavity linewidth $\delta\nu$ of this single mode depends on the FSR ($\nu_F \approx 902$\,MHz) and finesse ($F \approx 206$):

\begin{equation*}
  \delta\nu = \frac{\nu_F}{F} = \frac{902}{206}
\end{equation*}

\begin{equation*}
  \boxed{\delta\nu \approx 4.38\;\text{MHz}.}
\end{equation*}

% ============================================================================

\begin{question}[Question 7]Consider the same He-Ne laser of Problem 6 except that the longitudinal
modes of the laser may not coincide with the laser line. While the laser is running, the
frequencies of its longitudinal modes drift with time as a result of small thermally induced
changes in the length of the resonator. Find the allowable range of resonator lengths such that
the laser will always oscillate in one or two (but not more) longitudinal modes.
\end{question}

\bigskip
\textbf{\large Solution --- Question 7}

\medskip
\textbf{Allowable resonator lengths for 1 or 2 modes.}

Thermally induced drifts cause the modes to shift beneath the fixed atomic gain profile. To satisfy the constraint of "one or two (but not more)" modes, two extreme worst-case scenarios must be evaluated:

1. **Maximum Length (Preventing 3 modes):** The worst case for maximizing the number of oscillating modes is when one mode is centered perfectly at $\nu_o$. As established in Question 6, to prevent the adjacent modes ($\nu_o \pm \nu_F$) from reaching threshold, the cavity must not exceed:
\begin{equation*}
  L_{\max} \approx 16.6\;\text{cm}.
\end{equation*}

2. **Minimum Length (Guaranteeing at least 1 mode):** The worst case for minimizing the number of oscillating modes occurs when the modes straddle the peak symmetrically at $\nu_o \pm \nu_F/2$. To ensure the laser does not extinguish completely (dropping to zero modes), these straddling modes must have a gain equal to or greater than the loss:
\begin{equation*}
  \gamma_o\left(\nu_o \pm \frac{\nu_F}{2}\right) \ge \alpha_r.
\end{equation*}
\begin{equation*}
  2.5\times10^{-3} \exp\left[ -4\ln 2 \left( \frac{\nu_F/2}{\Delta\nu_D} \right)^2 \right] = \frac{0.01523}{L}.
\end{equation*}
Using $\frac{\nu_F/2}{\Delta\nu_D} = \frac{5}{L}$, this yields another transcendental equation:
\begin{equation*}
  2.5 \exp\left( -\frac{69.31}{L^2} \right) = \frac{15.23}{L}.
\end{equation*}
Solving numerically yields the minimum length required to guarantee oscillation despite thermal drift:
\begin{equation*}
  L_{\min} \approx 10.9\;\text{cm}.
\end{equation*}

Thus, the strict operational bounds are:

\begin{equation*}
  \boxed{10.9\;\text{cm} \le L \le 16.6\;\text{cm}.}
\end{equation*}

This elegantly illustrates the tight mechanical tolerances required to stabilize single-or-dual mode performance without active frequency locking.

% ============================================================================

\begin{question}[Question 8]An Ar$^{+}$-ion laser has a resonator of length 100\,cm. The refractive
index $n=1$.
\begin{enumerate}[label=(\alph*)]
  \item Determine the frequency spacing $\nu_{F}$ between the resonator modes.
  \item Determine the number of longitudinal modes that the laser can sustain if the FWHM
        Doppler-broadened linewidth is $\Delta\nu_{D}=3.5$\,GHz and the loss coefficient is
        half the peak small-signal gain coefficient.
  \item What would the resonator length $L$ have to be to achieve operation on a single
        longitudinal mode? What would that length be for a CO$_{2}$ laser that has a much
        smaller Doppler linewidth $\Delta\nu_{D}=60$\,MHz under the same conditions?
\end{enumerate}
\end{question}

\bigskip
\textbf{\large Solution --- Question 8}

\medskip
\textbf{Part (a): Mode spacing.}

\begin{equation*}
  \nu_F = \frac{c}{2L} = \frac{3\times10^{8}}{2(1)} = \boxed{150\;\text{MHz}.}
\end{equation*}

\textbf{Part (b): Number of oscillating modes.}

The loss coefficient is half the peak gain: $\alpha_r = \gamma_o(\nu_o)/2$. A mode at frequency $\nu$ oscillates when $\gamma_o(\nu) \ge \alpha_r = \gamma_o(\nu_o)/2$. For the Gaussian profile:

\begin{equation*}
  e^{-4\ln2\,(\nu-\nu_o)^2/\Delta\nu_D^2} \ge \tfrac{1}{2}
  \quad\Longrightarrow\quad |\nu-\nu_o| \le \frac{\Delta\nu_D}{2} = 1.75\;\text{GHz}.
\end{equation*}

The oscillation window has width $B = \Delta\nu_D = 3.5$\,GHz:

\begin{equation*}
  M = \frac{B}{\nu_F} = \frac{3.5\times10^{9}}{150\times10^{6}} \approx \boxed{23\;\text{modes}.}
\end{equation*}

\textbf{Part (c): Single-mode resonator lengths.}

For single-mode operation the mode spacing must exceed the oscillation window: $\nu_F > \Delta\nu_D$.

\begin{equation*}
  L_{\text{Ar}^+} < \frac{c}{2\Delta\nu_D} = \frac{3\times10^{8}}{2(3.5\times10^{9})} \approx \boxed{4.3\;\text{cm}.}
\end{equation*}

For CO$_2$ with $\Delta\nu_D = 60$\,MHz:

\begin{equation*}
  L_{\text{CO}_2} < \frac{3\times10^{8}}{2(60\times10^{6})} \approx \boxed{2.5\;\text{m}.}
\end{equation*}

% ============================================================================

\begin{question}[Question 9]A Doppler-broadened gas laser operates at 515\,nm in a resonator with two
mirrors separated by a distance of 50\,cm. The photon lifetime is 0.33\,ns. The spectral window
within which oscillation can occur is of width $B=1.5$\,GHz. To select a single mode, the light
is passed into an etalon (a passive Fabry-Perot resonator) whose mirrors are separated by the
distance $d$ and its finesse is $F$. The etalon acts as a filter. Suggest suitable values of $d$
and $F$. Is it better to place the etalon inside or outside the laser resonator?
\end{question}

\bigskip
\textbf{\large Solution --- Question 9}

\medskip
\textbf{Etalon design.}

The laser has FSR $\nu_F = c/2L = 300$\,MHz and oscillation window $B = 1.5$\,GHz ($M = 5$ modes). The etalon must satisfy two conditions simultaneously:

1. Its FSR must exceed the oscillation window so that only one laser mode fits within one etalon passband: $\nu_{F,e} = c/2d > B = 1.5$\,GHz:
\begin{equation*}
  d < \frac{c}{2B} = \frac{3\times10^{8}}{2(1.5\times10^{9})} = 10\;\text{cm}.
\end{equation*}
Choose $\boxed{d = 10\;\text{cm}}$.

2. The etalon linewidth must be narrower than the laser mode spacing $\nu_F = 300$\,MHz so it resolves individual modes: $\delta\nu_e = \nu_{F,e}/F_e < \nu_F$:
\begin{equation*}
  F_e > \frac{\nu_{F,e}}{\nu_F} = \frac{c/2d}{c/2L} = \frac{L}{d} = \frac{50}{10} = 5.
\end{equation*}
Choose $\boxed{F_e \approx 5\text{--}10}$ (a practical finesse of 10 provides a comfortable margin).

\textbf{Inside or outside?} The etalon is more effective \emph{inside} the resonator. When placed internally, a rejected mode experiences additional loss on every round trip, strongly suppressing it. Placed externally, it acts only as a passive filter after the light has already built up, requiring a much higher finesse to achieve the same mode-selection contrast.

% ============================================================================

\begin{question}[Question 10]Consider a 10-cm long gas laser operating at the centre of the 600-nm line
in a single longitudinal and a single transverse mode. The mirror reflectances are $R_{1}=99\%$
and $R_{2}=100\%$. The refractive index $n=1$ and the effective area of the output beam is
1\,mm$^{2}$. The small-signal gain coefficient $\gamma_{o}(\nu_{o})=0.1$\,cm$^{-1}$ and the
saturation photon-flux density $\Phi_{s}=1.43\times10^{19}$\,photons/cm$^{2}$-s.
\begin{enumerate}[label=(\alph*)]
  \item Determine the distributed loss coefficients $\alpha_{m1}$ and $\alpha_{m2}$ associated
        with each of the mirrors separately. Assuming that $\alpha_{s}=0$, find the resonator
        loss coefficient $\alpha_{r}$.
  \item Find the photon lifetime $\tau_{\mathrm{ph}}$.
  \item Determine the output photon flux density $\Phi_{\mathrm{out}}$ and the output
        power $P_{\mathrm{out}}$.
\end{enumerate}
\end{question}

\bigskip
\textbf{\large Solution --- Question 10}

\medskip
\textbf{Part (a): Mirror loss coefficients.}

Each mirror contributes a distributed loss coefficient $\alpha_{mi} = \ln(1/R_i)/(2L)$:

\begin{equation*}
  \alpha_{m1} = \frac{\ln(1/0.99)}{2(10)} = \frac{0.01005}{20} = 5.03\times10^{-4}\;\text{cm}^{-1},\qquad
  \alpha_{m2} = \frac{\ln(1/1.0)}{20} = 0.
\end{equation*}

\begin{equation*}
  \boxed{\alpha_r = \alpha_{m1} + \alpha_{m2} + \alpha_s = 5.03\times10^{-4}\;\text{cm}^{-1}.}
\end{equation*}

\textbf{Part (b): Photon lifetime.}

\begin{equation*}
  \tau_p = \frac{1}{c\,\alpha_r} = \frac{1}{(3\times10^{10})(5.03\times10^{-4})}
\end{equation*}

\begin{equation*}
  \boxed{\tau_p \approx 66.3\;\text{ns}.}
\end{equation*}

\textbf{Part (c): Output flux and power.}

At steady-state oscillation the saturated gain equals the loss:

\begin{equation*}
  \frac{\gamma_o}{1 + \Phi/\Phi_s} = \alpha_r
  \quad\Longrightarrow\quad
  \Phi = \Phi_s\left(\frac{\gamma_o}{\alpha_r} - 1\right)
  = (1.43\times10^{19})\left(\frac{0.1}{5.03\times10^{-4}} - 1\right)
  \approx 2.83\times10^{21}\;\text{photons/cm}^{2}\text{-s.}
\end{equation*}

The output photon flux density at mirror $M_1$ ($T_1 = 0.01$):

\begin{equation*}
  \Phi_{\text{out}} = T_1\,\Phi = 0.01(2.83\times10^{21}) = 2.83\times10^{19}\;\text{photons/cm}^{2}\text{-s.}
\end{equation*}

With photon energy $h\nu = hc/\lambda = 3.31\times10^{-19}$\,J and effective area $A = 1\;\text{mm}^2 = 10^{-2}\;\text{cm}^2$:

\begin{equation*}
  \boxed{P_{\text{out}} = h\nu\,\Phi_{\text{out}}\,A = (3.31\times10^{-19})(2.83\times10^{19})(10^{-2}) \approx 93.7\;\text{mW}.}
\end{equation*}

% ============================================================================

\begin{question}[Question 11]A He-Ne laser operating at 632.8\,nm produces 50\,mW of multimode power
at its output. It has an inhomogeneously broadened gain profile with a Doppler linewidth
$\Delta\nu_{D}=1.5$\,GHz. The resonator is 30\,cm long.
\begin{enumerate}[label=(\alph*)]
  \item If the maximum small-signal gain coefficient is twice the loss coefficient, determine
        the number of longitudinal modes of the laser.
  \item If the mirrors are adjusted to maximise the intensity of the strongest mode, estimate
        its power.
\end{enumerate}
\end{question}

\bigskip
\textbf{\large Solution --- Question 11}

\medskip
\textbf{Part (a): Number of longitudinal modes.}

With $L = 30$\,cm and $n=1$, the FSR is $\nu_F = c/2L = 500$\,MHz. The gain is twice the loss ($\gamma_o = 2\alpha_r$), so a mode at $\nu$ oscillates when $\gamma_o(\nu) \ge \alpha_r = \gamma_o/2$. Using the Gaussian profile (inhomogeneous broadening), this again gives an oscillation window of width $\Delta\nu_D$:

\begin{equation*}
  M = \frac{\Delta\nu_D}{\nu_F} = \frac{1.5\times10^{9}}{500\times10^{6}} \qquad\Longrightarrow\qquad \boxed{M = 3\;\text{modes}.}
\end{equation*}

\textbf{Part (b): Power of the strongest mode.}

For an inhomogeneously broadened medium each mode saturates independently. At steady state the intracavity flux of mode $q$ is $\Phi_q = \Phi_s(\gamma_o(\nu_q)/\alpha_r - 1)$. With $\nu_F/\Delta\nu_D = 500/1500 = 1/3$, the gain of the $\pm 1$ side modes relative to the centre mode is:

\begin{equation*}
  \frac{\gamma_o(\nu_o \pm \nu_F)}{\gamma_o(\nu_o)} = e^{-4\ln2/9} \approx 0.735
  \quad\Longrightarrow\quad \frac{\gamma_o(\nu_o \pm \nu_F)}{\alpha_r} = 2(0.735) = 1.469.
\end{equation*}

The flux ratios are $\Phi_0 : \Phi_{\pm1} = (2-1) : (1.469-1) = 1 : 0.469$, giving a relative power share for the centre mode of:

\begin{equation*}
  P_0 = \frac{1}{1 + 2(0.469)} \times 50\;\text{mW} = \frac{1}{1.938} \times 50\;\text{mW} \approx \boxed{25.8\;\text{mW}.}
\end{equation*}

% ============================================================================

\begin{question}[Question 12]Show that the effect of frequency pulling by the atomic medium is to
reduce the intermodal frequency separation from $c/2L$ to
$(c/2L)(1-\gamma c/2\pi\Delta\nu)$.
Calculate the reduction for the case of a laser with $\Delta\nu=10^{9}$\,Hz,
$\gamma=4\times10^{-2}$\,m$^{-1}$, and $L=100$\,cm.
\end{question}

\bigskip
\textbf{\large Solution --- Question 12}

\medskip
\textbf{Derivation.}

The round-trip resonance condition requires that the total phase accumulated over one round trip be a multiple of $2\pi$. With a dispersive medium of refractive index $n(\nu)$, the resonance frequencies satisfy $2L\,n(\nu_q)\,\nu_q/c = q$. Near resonance, the real part of the susceptibility $\chi'$ shifts the effective index: $n(\nu) \approx 1 + \chi'(\nu)/2$. For a Lorentzian line, $\chi'(\nu) \propto (\nu_o - \nu)/\Delta\nu$ with magnitude set by the gain coefficient $\gamma$. Differentiating the resonance condition with respect to mode number $q$ yields the corrected mode spacing:

\begin{equation*}
  \Delta\nu_{\text{modes}} = \frac{c}{2L}\left(1 - \frac{\gamma c}{2\pi\Delta\nu}\right).
\end{equation*}

\textbf{Numerical evaluation.}

The fractional reduction is $\gamma c/(2\pi\Delta\nu)$:

\begin{equation*}
  \frac{\gamma c}{2\pi\Delta\nu} = \frac{(4\times10^{-2})(3\times10^{8})}{2\pi(10^{9})} = \frac{1.2\times10^{7}}{6.283\times10^{9}} \approx 1.91\times10^{-3}.
\end{equation*}

The bare mode spacing is $c/2L = 150$\,MHz; the frequency pulling reduces it by:

\begin{equation*}
  \boxed{\delta(\Delta\nu) = \frac{c}{2L} \times 1.91\times10^{-3} \approx 150\;\text{MHz} \times 1.91\times10^{-3} \approx 286\;\text{kHz}.}
\end{equation*}

% ============================================================================

\begin{question}[Question 13]A laser cavity of length $L=10$\,cm and mirror reflectances $R_{1}=0.99$
and $R_{2}=0.85$ is filled with a gain medium of refractive index $n=3$ and scattering loss
factor $\alpha_{s}=10^{-4}$\,cm$^{-1}$.
\begin{enumerate}[label=(\alph*)]
  \item What is the FSR and $F$ of the cavity?
  \item Calculate its photon lifetime.
  \item Suppose that the active gain medium has a single-pass small-signal power gain $G_{o}$;
        how large should $G_{o}$ be to obtain laser oscillations?
\end{enumerate}
\end{question}

\bigskip
\textbf{\large Solution --- Question 13}

\medskip
\textbf{Part (a): FSR and finesse.}

With $n=3$ the FSR is:

\begin{equation*}
  \nu_F = \frac{c}{2nL} = \frac{3\times10^{10}}{2(3)(10)} = \boxed{5\times10^{8}\;\text{Hz} = 500\;\text{MHz}.}
\end{equation*}

The round-trip power loss factor is $\mathcal{L} = R_1 R_2\,e^{-2\alpha_s L} = (0.99)(0.85)e^{-2(10^{-4})(10)} \approx 0.8398$. The effective reflectance for the finesse formula is $\rho = \sqrt{\mathcal{L}} \approx 0.9164$:

\begin{equation*}
  F = \frac{\pi\rho^{1/2}}{1-\rho} \approx \frac{\pi(0.9574)}{1-0.9164} = \frac{3.008}{0.0836} \approx \boxed{36.}
\end{equation*}

\textbf{Part (b): Photon lifetime.}

The total round-trip loss exponent is $-\ln\mathcal{L} = -\ln(0.8398) = 0.1746$. The photon lifetime is:

\begin{equation*}
  \tau_p = \frac{2nL/c}{-\ln\mathcal{L}} = \frac{2(3)(10)/(3\times10^{10})}{0.1746} = \frac{2\times10^{-9}}{0.1746} \approx \boxed{11.5\;\text{ns}.}
\end{equation*}

\textbf{Part (c): Required single-pass gain.}

Oscillation occurs when the round-trip gain equals or exceeds all losses. After one round trip the field traverses the gain medium twice:

\begin{equation*}
  G_o^2 \cdot R_1 R_2 \cdot e^{-2\alpha_s L} \ge 1
  \quad\Longrightarrow\quad G_o \ge \frac{1}{\sqrt{R_1 R_2\,e^{-2\alpha_s L}}} = \frac{1}{\sqrt{0.8398}}
\end{equation*}

\begin{equation*}
  \boxed{G_o \ge 1.091 \quad(\text{a 9.1\% single-pass power gain}).}
\end{equation*}

% ============================================================================

\begin{question}[Question 14]Show that the gain factor of the amplified spontaneous emission (ASE),
$G_{o}-1$, has a bandwidth that is smaller than the bandwidth of the small-signal gain
coefficient $\gamma_{o}$, i.e.\ the amplification of spontaneous emission is accompanied by
spectral narrowing. You may assume that $\gamma_{o}(\nu)$ has a Lorentzian distribution.
\end{question}

\bigskip
\textbf{\large Solution --- Question 14}

\medskip
\textbf{Spectral narrowing of the ASE gain factor.}

For a Lorentzian gain coefficient:

\begin{equation*}
  \gamma_o(\nu) = \frac{\gamma_o(\nu_o)}{1+(2(\nu-\nu_o)/\Delta\nu)^2},
\end{equation*}

the FWHM of $\gamma_o(\nu)$ is exactly $\Delta\nu$. The ASE gain factor at each frequency is $G_o(\nu)-1 = e^{\gamma_o(\nu)L}-1$. Its half-maximum occurs where $e^{\gamma_{1/2}L}-1 = (G_o(\nu_o)-1)/2$, giving:

\begin{equation*}
  e^{\gamma_{1/2}L} = \frac{e^{\gamma_o(\nu_o)L}+1}{2}.
\end{equation*}

In the high-gain limit $G_o \gg 1$: $e^{\gamma_{1/2}L} \approx e^{\gamma_o L}/2$, so $\gamma_{1/2} = \gamma_o(\nu_o) - \ln2/L$. Since $\gamma_{1/2} > \gamma_o(\nu_o)/2$ (for any $G_o > 1$), the ASE half-maximum occurs \emph{closer} to line centre than the half-maximum of $\gamma_o(\nu)$ itself. Solving for the ASE bandwidth $\Delta\nu_{\text{ASE}}$ from the Lorentzian:

\begin{equation*}
  \Delta\nu_{\text{ASE}} = \Delta\nu\sqrt{\frac{\ln 2}{\gamma_o(\nu_o)L - \ln 2}} < \Delta\nu.
\end{equation*}

This confirms that \emph{ASE is spectrally narrowed relative to the atomic linewidth} --- a physical consequence of the exponential amplification preferentially boosting frequencies closest to the gain peak where $e^{\gamma_o L}$ is largest.

% ============================================================================

\begin{question}[Formative Question]\begin{enumerate}[label=(\alph*)]
  \item Write the rate equations for the two-level laser system shown in the figure.
  \item Show that a steady state population inversion cannot be achieved by using direct
        optical pumping between levels 1 and 2. Assume that level 1 is the ground state.
\end{enumerate}
\end{question}

\bigskip
\textbf{\large Solution --- Formative Question}

\medskip
\textbf{Part (a): Rate equations.}

Let $W = \sigma\Phi$ be the stimulated transition rate (equal for absorption and emission since levels have equal degeneracy), and $\tau$ be the upper-level spontaneous lifetime. With optical pumping between levels 1 (ground) and 2:

\begin{align*}
  \frac{dN_2}{dt} &= W N_1 - W N_2 - \frac{N_2}{\tau} = W(N_1 - N_2) - \frac{N_2}{\tau}\\[1ex]
  \frac{dN_1}{dt} &= -W N_1 + W N_2 + \frac{N_2}{\tau} = W(N_2 - N_1) + \frac{N_2}{\tau}
\end{align*}

with the conservation constraint $N_1 + N_2 = N_a$.

\textbf{Part (b): Impossibility of steady-state inversion.}

At steady state ($dN_2/dt = 0$):

\begin{equation*}
  W(N_1 - N_2) = \frac{N_2}{\tau}
  \quad\Longrightarrow\quad
  \frac{N_2}{N_1} = \frac{W\tau}{1 + W\tau} < 1.
\end{equation*}

Since $W\tau/(1+W\tau) < 1$ for any finite $W$, we always have $N_2 < N_1$. Even in the limit $W\tau \to \infty$ (infinite pump intensity), $N_2/N_1 \to 1$, meaning the populations merely equalise but never invert. This result is fundamental: a two-level system cannot achieve population inversion by direct optical pumping because the pump simultaneously stimulates emission from the very upper level it is trying to fill, creating a self-limiting saturation ceiling at $\Delta N = 0$. A minimum of three levels is required to decouple the pump and the lasing transitions.
